% The preamble: packages, page layout, draft notes, theorem environments, cross-reference names and generic letters.
% Loaded by the main file before the notation. Only the packages loaded here are available to the sections.

\usepackage[margin=0.9in]{geometry}
\usepackage[utf8]{inputenc}
\usepackage[T1]{fontenc}

%-------------------------------------------------------------------------------
% Packages
\usepackage[dvipsnames]{xcolor}
\usepackage{amssymb}
\usepackage[mathscr]{eucal}
\usepackage{amsthm}
\usepackage{amsmath}
\usepackage{mathtools}
\definecolor{highlight}{HTML}{3465a4}
\usepackage[hypertexnames=false, hyperfootnotes=false, colorlinks=true, linktocpage=true, allcolors=highlight, bookmarksdepth=2]{hyperref}
\usepackage{cleveref}
\usepackage{url}
\usepackage{array}

%-------------------------------------------------------------------------------
% Typefaces and layout
\usepackage{mathpazo}
\usepackage{courier}
\usepackage{microtype}
\linespread{1.06}
\setlength{\parskip}{4pt}
\setlength{\parindent}{0pt}
\setcounter{secnumdepth}{3}
% enumitem is not assumed: first-level enumerate is labelled (i), (ii), ...; for named parts use
% description or \item[(a)] with an explicit label.
\renewcommand{\labelenumi}{(\roman{enumi})}
\renewcommand{\theenumi}{(\roman{enumi})}
\numberwithin{equation}{section}

%-------------------------------------------------------------------------------
% Draft notes: visible, coloured, usable in text and in math. Set \showdraftnotesfalse to suppress them.
\newif\ifshowdraftnotes
\showdraftnotestrue
\newcommand{\draftnotetext}[1]{\textcolor{red}{\footnotesize[\textbf{Draft note:} #1]}}
\newcommand{\draftnote}[1]{\ifshowdraftnotes\ifmmode\text{\draftnotetext{#1}}\else\draftnotetext{#1}\fi\fi}

%-------------------------------------------------------------------------------
% Theorem environments
\theoremstyle{plain}
\newtheorem{thm}{Theorem}[section]
\newtheorem{lem}[thm]{Lemma}
\newtheorem{prop}[thm]{Proposition}
\newtheorem{cor}[thm]{Corollary}
\theoremstyle{definition}
\newtheorem{defn}[thm]{Definition}
\newtheorem{rmk}[thm]{Remark}
\newtheorem{notation}[thm]{Notation}
\newtheorem{example}[thm]{Example}
\newtheorem{conv}[thm]{Convention}

\crefname{equation}{Eq.}{Eqs.}
\crefname{thm}{Theorem}{Theorems}
\crefname{lem}{Lemma}{Lemmas}
\crefname{prop}{Proposition}{Propositions}
\crefname{cor}{Corollary}{Corollaries}
\crefname{defn}{Definition}{Definitions}
\crefname{rmk}{Remark}{Remarks}
\crefname{notation}{Notation}{Notations}
\crefname{example}{Example}{Examples}
\crefname{conv}{Convention}{Conventions}
\crefname{section}{Section}{Sections}
\crefname{subsection}{Section}{Sections}
\crefname{appendix}{Appendix}{Appendices}
\crefname{table}{Table}{Tables}

%-------------------------------------------------------------------------------
% Generic letters
\newcommand{\bbA}{\mathbb{A}}
\newcommand{\bbC}{\mathbb{C}}
\newcommand{\bbD}{\mathbb{D}}
\newcommand{\bbE}{\mathbb{E}}
\newcommand{\bbF}{\mathbb{F}}
\newcommand{\bbK}{\mathbb{K}}
\newcommand{\bbL}{\mathbb{L}}
\newcommand{\bbP}{\mathbb{P}}
\newcommand{\bbQ}{\mathbb{Q}}
\newcommand{\bbR}{\mathbb{R}}
\newcommand{\bbT}{\mathbb{T}}
\newcommand{\bbZ}{\mathbb{Z}}
\newcommand{\cA}{\mathcal{A}}
\newcommand{\cC}{\mathcal{C}}
\newcommand{\cE}{\mathcal{E}}
\newcommand{\cF}{\mathcal{F}}
\newcommand{\cG}{\mathcal{G}}
\newcommand{\cL}{\mathcal{L}}
\newcommand{\cN}{\mathcal{N}}
\newcommand{\cO}{\mathcal{O}}
\newcommand{\cP}{\mathcal{P}}
\newcommand{\cQ}{\mathcal{Q}}
\newcommand{\cW}{\mathcal{W}}
\newcommand{\cZ}{\mathcal{Z}}
\newcommand{\mgp}[1]{\mathsf{#1}}

%-------------------------------------------------------------------------------
% General
\newcommand{\dd}{\mathrm{d}}
\newcommand{\ri}{\mathrm{i}}
\newcommand{\re}{\mathrm{e}}
\newcommand{\de}{\partial}
% \llbracket: from a font package if available, otherwise a substitute (stmaryrd is not assumed)
\providecommand{\llbracket}{[\![}
\providecommand{\rrbracket}{]\!]}
\newcommand{\fps}[1]{\llbracket #1 \rrbracket}
\newcommand{\lser}[1]{(\!(#1)\!)}
\newcommand{\Frac}{\mathrm{Frac}}
\DeclareMathOperator{\Tr}{Tr}
\DeclareMathOperator{\tr}{tr}
\DeclareMathOperator{\End}{End}
\DeclareMathOperator{\Aut}{Aut}
\DeclareMathOperator*{\Res}{Res}
\DeclareMathOperator{\val}{val}
\DeclareMathOperator{\ev}{ev}
\DeclareMathOperator{\diag}{diag}
\DeclareMathOperator{\ord}{ord}
\newcommand{\tensor}{\otimes}
\newcommand{\coef}[1]{[#1]}                    % coefficient extraction, e.g. \coef{z^n}
